\documentclass[11pt]{article}
\usepackage{fontspec,amsmath,amssymb,amsthm,hyperref}
\usepackage[margin=1in]{geometry}
\setlength{\emergencystretch}{3em}
\title{Exact KIPLMC1 clock change for correlated Gaussian innovations}
\author{LangevinLean mathematical exposition}
\date{}
\begin{document}\maketitle
\noindent Source module: \texttt{KIPLMC/ExponentialRescaling.lean}. Compile this proof document with XeLaTeX or LuaLaTeX. The accompanying HTML pairs the mathematical argument with the exact Lean statements and records the scope of the formalization.


The same positions and a velocity multiplication by √N convert physical-time KIPLMC1 into the shared unit-temperature exponential integrator for the sum potential. This file proves the actual covariance, Gaussian measure, update, initialization, and complete iterated-law identities. Both correlated noise components and their mutual covariance are transported explicitly.


\section{The Brownian clock changes each correlated block diagonally}


Write \(c>0\), \(\gamma^\prime=c\gamma\), and \(h^\prime=h/c\). A time change contributes a Brownian factor \(c^{-1/2}\); the position response contributes another factor \(c^{-1}\). The resulting map on the two-component Gaussian block is \[D_c(v,x)=(c^{-1/2}v,c^{-3/2}x).\] Lean constructs this diagonal map as an actual continuous linear map on Euclidean two-space. The coordinate theorem verifies the factors directly. The kernel accessor theorem identifies the two entries of the algorithmic kernel vector with \(\psi_0\) and \(\psi_1\), so the subsequent proof can use the already checked arbitrary-index response scaling.



\noindent Lean declarations: \texttt{kineticBlockScale}, \texttt{kineticBlockScale\_apply}, \texttt{klmcKernels\_apply}, \texttt{phaseClockScale}.

\section{Prove the covariance transformation and the Gaussian measure identity}


The covariance is the actual Gram integral \(C_{ij}(\gamma,h)=\int_0^h\psi_i(\gamma,t)\psi_j(\gamma,t)\,dt\), with \(i,j\in\{0,1\}\). The exact response identity \(\psi_i(c\gamma,t)=c^{-i}\psi_i(\gamma,ct)\), followed by the substitution \(s=ct\), gives \[C_{ij}(c\gamma,h/c)=c^{-(i+j+1)}C_{ij}(\gamma,h).\] These are precisely the entries of \(D_c C D_c^\mathsf T\), including the off-diagonal velocity-position covariance. Lean rewrites both response factors, moves their constants outside the integral, performs the checked interval-integral change of variables, and reduces the four finite index cases using \((\sqrt c)^2=c\).



Both measures are actual centered Gaussian measures. The mean of their linear pushforward is zero by commutation of the Bochner integral with the continuous linear map. Their covariance bilinear forms agree on the standard orthonormal basis by the displayed Gram identity, so Gaussian uniqueness proves \[(D_c)_\#\mathcal G_{\gamma,h}=\mathcal G_{c\gamma,h/c}.\] No Gaussian covariance or distributional scaling is introduced as a premise.



\noindent Lean declarations: \texttt{klmcCovariance\_clock}, \texttt{klmcGaussian\_clock}.

\section{Transport the independent spatial blocks without losing correlation}


The particle noise measure is the finite product of these two-component Gaussian blocks over the parameter and latent coordinate set. Lean applies the proved scalar Gaussian pushforward to each independent block, uses the exact product-measure mapping identity, and checks that assembling blocks into velocity and position vectors commutes with \(D_c\). Thus the full actual particle innovation obeys the same clock law. Independence between distinct spatial blocks is retained, while each block retains its exact velocity-position correlation.



\noindent Lean declarations: \texttt{exponentialNoise\_clock}.

\section{Conjugate the actual exponential update}


For \(c=\sqrt N\), multiply velocity by \(c\) and leave position fixed. Under the changed clock, the force becomes \(\nabla V_N\), the friction is \(c\gamma\), and the step is \(h/c\). The exact deterministic exponential response conjugacy has already been proved. What remains is the actual additive innovation: \[\sqrt{2c\gamma}\,D_c\xi=\bigl(c\sqrt{2\gamma/N}\,\xi_v,\sqrt{2\gamma/N}\,\xi_x\bigr).\] Lean proves these two scalar identities from the positive square roots and the checked particle amplitude identity, then rewrites the shared exponential map. The resulting equality is an identity of the implemented updates for every phase state and every supplied correlated innovation.



\noindent Lean declarations: \texttt{kiplmc1Step\_unitTemperature}.

\section{Pass the clock identity through the actual Markov laws}


The generic innovation-mapping lemma proves that replacing the noise law by \(R_\#Q\) is exactly substituting \(R\xi\) in each transition. Its proof is an induction on the true iterated Markov measure: at each step, the product law is pushed through \((z,\xi)\mapsto(z,R\xi)\), then through the measurable update. The phase conjugacy theorem applies this fact together with the proved update identity and the exact Gaussian block scaling. Consequently the entire actual KIPLMC1 Gaussian Markov law, pushed through \((v,x)\mapsto(\sqrt N v,x)\), is exactly the shared unit-temperature KLMC Markov law for \(V_N\) at \((\sqrt N\gamma,h/\sqrt N)\). The equality includes the noise measure and every iteration, rather than merely comparing update formulas.



\noindent Lean declarations: \texttt{markovLaws\_noise\_map}, \texttt{kiplmc1GaussianLaws\_unitTemperature}.

\section{Identify the thermal initialization and the unchanged position marginal}


The physical velocity target is the explicit Gaussian pushforward \(N^{-1/2}\mathcal N(0,I)\). Multiplication by \(\sqrt N\) returns the actual standard Gaussian law. Since the position is left fixed, the independent thermal initialization becomes \(\mathcal N(0,I)\otimes\nu\) for every initial position law \(\nu\). Lean proves this using the actual product-measure pushforward and the already checked thermal velocity normalization. Finally, applying the position projection to the full iterated-law identity proves that the actual KIPLMC1 position law equals the unit-temperature KLMC position law at the changed clock. This supplies the exact law transfer needed to import the shared kinetic error estimate without duplicating its analysis.



\noindent Lean declarations: \texttt{thermalInitial\_unitTemperature}, \texttt{kiplmc1GaussianPositionLaws\_unitTemperature}.

\end{document}
